\section{Limitations of the Prototype}
\label{sec:limitations-of-prototype}

\minisec{Limited functional scope}
The prototype is deliberately limited in scope and, as a minimum viable product, is aimed primarily
at hidden profiles, so that the focus and the subsequent validation remain attributable
(cf.\ \autoref{sec:hidden-profiles}).
This scoping assumes that the phenomenon can be examined in isolation. That assumption
is presumably too simple: in real meetings, shared-information bias overlaps with status and
hierarchy effects, confirmation bias and groupthink, as well as time pressure, drifting off topic, and
an unclear goal (cf.\ \autoref{sec:meetings-in-modern-work}).
The present version covers this interplay, if at all, only implicitly.
In addition, as generally available AI systems have advanced, people have grown used to speaking
freely with language models and to expecting answers to any concern.
That expectation can exceed the narrow repertoire of the prototype. Participants
might assume response and interaction options that the system does not provide, and thus, because of
this perceived gap, lose trust even in existing functions of the
system that are not actually deficient.

\minisec{Unimodal perception}
The system works exclusively on the transcript.
Transcription errors, for instance from unclear pronunciation or poor microphone quality, enter
directly into the context of the LLM pipeline and can distort thoughts and interventions.
Prosody (stress, pronunciation, tone of voice) and visual signals such as facial expression, gaze, and
visible attention or lack of interest are absent.
Facilitators, however, use precisely these signals to read how an utterance is meant and whether
someone is present (cf.\ \autoref{ch:empirical-analysis-of-facilitating-strategies}).
A sentiment analysis of the audio stream was not implemented.
This is an interpretive gap. It is, however, in line with the concerns about control and
surveillance expressed in the interviews.

\minisec{Runtime mechanics}
The thought pipeline is multi-stage and probabilistic. Interventions can therefore arrive
late, and the same conversational situation can turn out differently.
FIM scores are assigned once and then only aged deterministically. They are not
re-evaluated semantically (cf.\ \autoref{sec:dynamic-im}).

\minisec{Speaking over the agent}
If a participant speaks into an ongoing utterance of the agent, the system
does not break it off.
The thought already selected is spoken through until the next human turn.
Only this new element starts the pipeline again, so that a further thought can be formed and
uttered.
An overlapping turn therefore does not change the contribution in progress, and it can lead to odd or
unnatural interactions when participants continue speaking while the AI facilitator is speaking.

\minisec{Custom meeting platform}
The meeting runs in a purpose-built web application and not on established meeting platforms such as
Zoom or Teams. This secures control over interaction and UX. Participants, however, are not used to
this new application, unlike the familiar tools.
Comparability with the familiar situation of use is therefore limited.

\minisec{Modified turn-allocation prediction}
The prediction of who speaks next is based on Liu et al.\
\cite{liuProactiveConversationalAgents2025}.
The prompt, however, had to be adapted from everyday dialogue to the facilitation context.
It therefore cannot be assumed that the results of the paper transfer one-to-one to this context.
A closer statistical validation would be advisable, but it was not possible for lack of a suitable
domain-specific corpus of facilitation settings with turn-allocation notation.

\minisec{Assumption of equal speaking shares}
The participation scan assumes that a balanced distribution of speaking time is
beneficial to the meeting (cf.\ \autoref{sec:participation-scan}).
This assumption holds for the decision format under evaluation, in which unshared information
enters the discussion only through spoken turns.
It does not hold for formats whose purpose is unequal speaking time, such as presentations or
longer contributions from experts.
There, curtailment and invitation would trigger even though the unequal distribution is intended.

\minisec{Limited evidence that the thresholds and prompts are ideal}
The prompts are based on the original prompts of the Inner Thoughts framework and were adapted
intuitively on the basis of the findings from the literature and the expert interviews.
Some uncertainty nevertheless remains as to whether these prompts were chosen ideally and produce
the effect they are meant to have. This is partly due to the limitations of context engineering
(cf.\ \autoref{sec:context-engineering}).
In addition, the adaptation of the prompts is shaped by the author and was adjusted through the
EDD procedure.
